\section{Fixed tails of size at least seven with arbitrary tensor partners}
\label{sec:factorwise-small-tail}

\paragraph{Status and exact scope.}
HD-FT04--06, reserved question Q2de, is a new derivation by
\texttt{/root/g34\_eightfolds}.  The coefficient-polynomial input is the
already checked D-AFT01, and the actual triple-cover genus argument is
the root's HD-FT01.  The new estimate controls the number of inertia
entries at which the polynomial comparison has an error; it does not
transfer twelve branch marks to a tensor factor.
Complete author and separate root whole-proof readings pass. Root also
read the complete source/status record, the entire supplementary checker
and its first successful output; no fresh root execution is claimed.

Use exactly the actual quotient and factorization of (FT.1)--(FT.2):
a connected finite-group cover of the sphere, an actual quotient onto
$P=A_n$ or $S_n$, an actual coefficient factorization $V=W\otimes U$
up to the permitted scalar factors, and genuine inertia $T_i$ with
\[
 e=\dim V,\qquad
             \sum_i\operatorname{rank}(1-T_i)=2e+6.
 \tag{FTL.1}
\]
The tensor partner $U$, its degree and coefficient field are unrestricted.
No primitivity, factorwise tensor-indecomposability, transitive equal-factor
core, or twelve-mark condition on the projected permutations is assumed.
An abstract spectrum without the actual quotient cover does not suffice.

Fix $\mu\vdash k$, $k\ge7$, and let $W=S^{(n-k,\mu)}$, or its
irreducible ordinary restriction to $A_n$.  Transposing the full partition
only twists by sign on $S_n$, which changes none of the absolute-character
or projective-support bounds below.  The threshold will imply $n>2k$,
so the indicated partition and its irreducible alternating restriction
are in the stable, non-self-conjugate range.

\subsection*{HD-FT04. A localized cubic character bound}

Write the exact Garsia--Goupil polynomial, in the notation of D-AFT01, as
\[
 \chi_W(x)=q_\mu(X_1(x),\ldots,X_k(x)),\quad
 q_\mu=C X_1^k+R,\quad C=f^\mu/k!,\quad
 L=\|R\|_1,\quad M=1+L/C.
 \tag{FTL.2}
\]
Here $\|R\|_1$ is the sum of the absolute values of its ordinary-monomial
coefficients.  The published formula and its proof give weighted degree
at most $k$, with weight $i$ on $X_i$; the sole ordinary-degree-$k$
monomial is $C X_1^k$.  Thus $R$ has ordinary degree at most $k-1$.
These are precisely the coefficient premises already proved in D-AFT01,
using \cite[Proposition~I.1 and its complete proof, pp.~9--11]{HD:GarsiaGoupil2009}.

\begin{lemma}[HD-FT04]
Suppose $k\ge7$ and
\[
                         n\ge2^{k+1}M.
 \tag{FTL.3}
\]
Put $a=\dim W=q_\mu(n,0,\ldots,0)$ and
$\varepsilon=2M/n^4$.  Then $a\ge2n^3$, and for every permutation
$x$, writing $f=X_1(x)$,
\[
 \frac{|\chi_W(x)|}{a}\le
 \begin{cases}
    (f/n)^3,& f\ge n/2,\\
    (f/n)^3+\varepsilon,& f<n/2.
 \end{cases}
 \tag{FTL.4}
\]
These bounds include the identity and all powers of every permutation.
\end{lemma}

\begin{proof}
Set $t=f/n$ and $m=n-f$.  All cycle counts are at most $n$.
The degree premise gives
\[
 |R(n,0,\ldots,0)|\le Ln^{k-1},\qquad
        a\ge Cn^k(1-M/n)\ge\tfrac12 Cn^k.
 \tag{FTL.5}
\]
We used $n\ge2^{k+1}M\ge256M$.  Also $C\ge1/k!$ and
$k!\le2^{k(k-1)/2}$, since $j\le2^{j-1}$ for each $j\ge1$.
Consequently
\[
 n^{k-3}\ge2^{(k+1)(k-3)}\ge4k!,
 \qquad
 (k+1)(k-3)-2-k(k-1)/2=(k-5)(k+2)/2>0.
 \tag{FTL.6}
\]
Equations (FTL.5)--(FTL.6) prove $a\ge2n^3$.
They also give $n>2k$, so all stable-character premises hold.

Suppose first that $f\ge n/2$.  The crude remainder bound gives
\[
 q_\mu(f,X_2,\ldots)\ge Cf^k-Ln^{k-1}>0,
 \tag{FTL.7}
\]
because $n/2^k\ge2M>L/C$.
For a pure monomial $f^j$ of $R$,
$|f^j-n^j|\le jmn^{j-1}\le k m n^{k-2}$; its constant term
has zero difference.  Every mixed monomial has ordinary degree at
most $k-1$ and contains a non-$f$ cycle count at most $m$, so its
value is at most $m n^{k-2}$ and its value at $(n,0,\ldots)$ is zero.
Summing coefficientwise yields
\[
 |R(f,X_2,\ldots)-R(n,0,\ldots)|\le k Lm n^{k-2}.
 \tag{FTL.8}
\]
Since $1-t^3\le3m/n$, the difference between the desired character
bound and its leading term is controlled as follows:
\[
 \begin{split}
 q_\mu-a t^3
 &\le-Cn^k(t^3-t^k)+(k+3)Lm n^{k-2}\\
 &\le-Cm n^{k-1}/8+(k+3)Lm n^{k-2}\le0.
 \end{split}
 \tag{FTL.9}
\]
Indeed $t\ge1/2$ and $k\ge4$ give
$t^3-t^k=t^3(1-t^{k-3})\ge(1-t)/8$.
Moreover $2^{k+1}\ge8(k+3)$ for $k\ge7$, so (FTL.3) implies
$n\ge8(k+3)L/C$.  When $m=0$ the difference is exactly zero.
Together with (FTL.7), this proves the first line of (FTL.4),
including its absolute value.

Now assume $0\le t<1/2$.  Consider any monomial of $R$ and let $j$
be its exponent of $f$, and $b$ the sum of its other exponents.
Then $j+2b\le k$ and $j+b\le k-1$.
After division by $n^k$, its absolute value is at most
$t^j n^{j+b-k}$.  If $j\ge3$, this is at most $t^3/n$.
If $j=0$, it is at most $n^{-\lceil k/2\rceil}\le n^{-4}$.
If $j=1$ or $2$, it is at most $t/n^3$ or $t^2/n^3$, respectively.
The following exact nonnegative identities control those two cases:
\[
 \begin{split}
 \frac{t^3}{n}+\frac2{n^4}-\frac{3t}{n^3}
     &=\frac{(nt-1)^2(nt+2)}{n^4}\ge0,\\
 \frac{2t^3}{n}+\frac1{n^7}-\frac{3t^2}{n^3}
     &=\frac{(n^2t-1)^2(2n^2t+1)}{n^7}\ge0.
 \end{split}
 \tag{FTL.10}
\]
Thus every normalized remainder monomial is at most
$t^3/n+n^{-4}$, and coefficientwise summation gives
\[
 \frac{|q_\mu|}{a}
 \le\frac{t^k+(M/n)t^3+M/n^4}{1-M/n}
 \le t^3+\frac{2M}{n^4}.
 \tag{FTL.11}
\]
For the last inequality, $t^k\le t^3/16$ for $k\ge7$,
and $M/n\le1/4$.  Hence the coefficient of $t^3$ is at most
$(1/16+1/4)/(3/4)=5/12<1$, while the last coefficient is at most
$4M/(3n^4)\le2M/n^4$.  This proves the second line of (FTL.4).
The argument depends only on cycle counts, so it applies to every power.
\end{proof}

\subsection*{HD-FT05. At most ten error-bearing inertia entries}

\begin{theorem}[HD-FT05]
Under the actual quotient and factorization (FTL.1), the factor $W$
described above is impossible whenever
\[
                 n>3500000,\qquad n\ge2^{k+1}M,
                 \qquad k\ge7.
 \tag{FTL.12}
\]
The conclusion allows every original imaginary quadratic field and
arbitrary tensor partners and scalar twists.
\end{theorem}

\begin{proof}
Let $x_i\in P$ be the actual projected inertia.  Call an index $i$
exceptional for this proof if $f(x_i)<n/2$, and let $b$ count these
indices.  This terminology records only where the error in (FTL.4)
can appear; it asserts no new conjugacy classification.
If an index is not exceptional, every power of $x_i$ still fixes
at least $n/2$ letters, so every power has zero error in (FTL.4).
At an exceptional index the Fourier average has error at most
$\varepsilon=2M/n^4$.

For any $x$ of finite order $h$, cyclic Fourier inversion bounds
the multiplicity of every target eigenvalue by
$h^{-1}\sum_{j=0}^{h-1}|\chi_W(x^j)|$.  Since $0\le f(x^j)/n\le1$
and its cube is at most itself, HD-FT04 gives at every exceptional index
\[
 \frac{\nu(W(x))}{a}
   \ge1-\frac1h\sum_j\frac{f(x^j)}n-\varepsilon
   =\frac{\operatorname{ind}(x)}n-\varepsilon
   >\frac14-\varepsilon>\frac15.
 \tag{FTL.13}
\]
The equality is the natural orbit-counting formula.  A permutation with
$n-f$ moving letters has index at least $(n-f)/2$; the exceptional
inequality is strict because $f<n/2$.
For the final strict inequality note that $n\ge256M$, so
\[
        \varepsilon\le\frac1{128n^3}<\frac1{20},
        \qquad
        2+\frac6e\le2+\frac3{n^3}<\frac{11}{5}.
 \tag{FTL.14}
\]
Here $e\ge a\ge2n^3$ by HD-FT04, and $n\ge256$.
Spectral convolution for the actual tensor factorization gives
$\operatorname{rank}(1-T_i)/e\ge\nu(W(x_i))/a$ for every $i$,
with any scalar twists.  Eleven exceptional entries would therefore
make (FTL.1) strictly greater than $11e/5$, contrary to (FTL.14).
Thus
\[
                              b\le10.
 \tag{FTL.15}
\]
No lower bound on the number of genuine projected marks was used.

Let $\Omega=\{1,\ldots,n\}^3$, with repetitions allowed.
The exact orbit-counting formula and (FTL.4), now retaining the cubes,
give for each inertia
\[
 \frac{\operatorname{rank}(1-T_i)}e
 \ge\frac{\nu(W(x_i))}{a}
 \ge\frac{\operatorname{ind}(x_i,\Omega)}{n^3}
                  -\varepsilon\,\mathbf1_{f(x_i)<n/2}.
 \tag{FTL.16}
\]
For $x_i=1$ both terms on the right are zero, so no issue is caused
by original marks disappearing in the quotient.

The geometric step of HD-FT01 applies to this same actual quotient.
There are five $P$-orbits on $\Omega$, of genera
$\gamma_1,\gamma_2,\gamma_2,\gamma_2,\gamma_3$.
The unordered-three-subset quotient has genus at least two for
$n>3500000$, by the published Monderer--Neftin theorem in exactly
the source scope recorded in HD-FT01.\cite[Theorem~1.2 and
Remark~2.8]{HD:MondererNeftin2022}
Indeed its stabilizer has order at most $6(n-3)!<|A_{n-2}|$ and
cannot contain either required alternating subgroup.
The actual forget-order morphism has degree six for $A_n$ as well as
$S_n$, so $\gamma_3\ge7$ by Riemann--Hurwitz.  The five actual
quotient covers therefore satisfy
\[
 \sum_i\operatorname{ind}(x_i,\Omega)
       =2(n^3-5)+2(\gamma_1+3\gamma_2+\gamma_3)
       \ge2n^3+4.
 \tag{FTL.17}
\]
This is a quotient-cover statement, not an assertion about a formal
list of permutations.

Summing (FTL.16) and using (FTL.15)--(FTL.17) yields
\[
 \frac1e\sum_i\operatorname{rank}(1-T_i)
 \ge2+\frac4{n^3}-\frac{20M}{n^4}
 >2+\frac3{n^3}\ge2+\frac6e.
 \tag{FTL.18}
\]
The strict inequality follows from $n\ge256M>20M$.
This contradicts (FTL.1) and proves the theorem.
\end{proof}

\subsection*{HD-FT06. The remaining unbounded ordinary tails have size at most six}

\begin{corollary}[HD-FT06]
In the actual factorwise tensor setting, suppose an ordinary
symmetric or alternating irreducible factor has natural degree
\[
                              n\ge241^{482}.
 \tag{FTL.19}
\]
Then its surviving partition label is extendible, and must satisfy
\[
                  \min\{n-\lambda_1,n-\lambda'_1\}\le6.
 \tag{FTL.20}
\]
The conclusion includes arbitrary factorwise partners, unequal
partner dimensions and all scalars.  It does not assert that any
of the remaining small-tail sizes occurs.
\end{corollary}

\begin{proof}
HD-FT02 removes all extendible labels with both tails at least241,
and HD-FT03 removes the split alternating constituents already for
$n>3500000$.  It remains to consider a label, transposed if necessary,
of the form $(n-k,\mu)$ with $7\le k\le240$.
The coefficient estimate (TW.23)--(TW.24) gives the uniform bound
\[
 M\le1+k!(k+1)2^{k-1}
        \le M_*:=1+241\cdot240!\cdot2^{239}.
 \tag{FTL.21}
\]
This is the accepted fixed-subset/branching calculation and does not
require enumerating all partitions of 240.
The following elementary comparison suffices:
\[
 2^{k+1}M\le2^{241}M_*
       \le2^{481}241^{241}<241^{482}.
 \tag{FTL.22}
\]
For the middle inequality use $M_*=1+241!2^{239}\le2^{240}241!$
and $241!\le241^{241}$.  For the strict final inequality use
$241^{241}>4^{241}=2^{482}>2^{481}$.
Also $241^{482}>3500000$.  Hence every such label satisfies
(FTL.12), and HD-FT05 excludes it.  This leaves exactly the necessary
small-tail range (FTL.20).
\end{proof}

\paragraph{Boundary and provenance.}
The new ingredient is the character-polynomial comparison with an error
supported on at most ten actual inertia entries.  The proof uses the
already accepted character-polynomial formula and quotient-genus theorem,
not an effective constant inferred from an asymptotic genus estimate.
It removes the previously unbounded fixed tails seven through 240
in the factorwise tensor-decomposable gate.  Tails zero through six,
other projective factor types and the actual quotient/factorization
interface remain separate.  No factor is assumed to be a sixfold
source on its own, and no independent positive cycle is constructed.
